\section{The first NC step for the actual ambient-dimensional canonical invariant}
\label{sec:hironaka-normal-crossings-canonical-initial-step}

Let $H$ be an actual NC carrier in a complex manifold of dimension $n$,
and let $M_H=(X,X,\varnothing,I_H,1)$ be its unchanged original marked datum.
For a proved resolution invariant $R$ whose dimension bound is $k\geq n$,
write $V^R_H(x)=R(M_H,\mathrm{refl},x)$. Its value space is arbitrary.
The ambient-dimensional canonical invariant is specifically
$R_n=\mathrm{canonicalInvariant}(n)$, with values in $\mathrm{ValTy}(n)$.
It is distinct from the previously considered family of depth $n+1$;
no equality between invariant depths is assumed. All histories below are
the initial reflexive history or the actual single step constructed from it.

\begin{lemma}[Whole NC chart transport for any proved invariant]
\label{hir:general-nc-generic-initial-value-flat-chart}
\leanok
\lean{Hironaka.HironakaNormalCrossingsCanonicalInitialStep.initialValue_eq_coordinate_model}
Suppose an actual centred flattening at $x\in H$ identifies the whole
carrier on its open chart domain with the coordinate union $H_S$.
For every proved resolution invariant $R$ of dimension bound $k\geq n$,
$V^R_H(x)=V^R_{H_S}(0)$, with its entire value retained.
\uses{hir:general-nc-standard-coordinate-carrier,hir:general-nc-initial-full-value-flat-chart}
\end{lemma}
\begin{proof}
The actual NC original-domain theorem gives a well-formed marking-one
datum of dimension at most $n$. Increasing the dimension bound to $k$
preserves each domain condition. The actual chart and its actual inverse
give a biholomorphism from the open flat domain to the lifted open target,
taking $x$ to zero and the whole carrier trace to the trace of $H_S$.
Composition with the chart and its inverse identifies the whole vanishing
ideal, as a predicate on holomorphic germs. The pulled-back marked data
therefore agree in all four fields: full ambient carrier, empty boundary,
whole ideal and marking one. The analogous whole-record identity under
an open inclusion follows by restricting the same vanishing predicate.

The reflexive tower is canonical for $R$ without an assumption on values.
Apply $R$'s proved restriction law to the original open domain, its whole
tower-isomorphism law to the actual biholomorphism, and its restriction
law to the coordinate-model target. The points involved belong to their
literal original cosupports, which equal the corresponding NC carriers.
The three whole-value equalities and the whole-record identities give
the asserted equality. No formula for a value or equality of two different
invariant families is used.
\end{proof}

\begin{lemma}[Whole highest-stratum constancy for any proved invariant]
\label{hir:general-nc-generic-initial-highest-stratum-constant}
\leanok
\lean{Hironaka.HironakaNormalCrossingsCanonicalInitialStep.initialValue_eq_of_highest_branch_count}
Suppose $r>0$ bounds every original branch count. Any two points of count
$r$ have the same entire initial value for every proved invariant $R$
of dimension bound at least $n$.
\uses{hir:general-nc-generic-initial-value-flat-chart,hir:general-nc-highest-stratum-branch-adapted-coordinates}
\end{lemma}
\begin{proof}
At each of the two points, the actual highest-stratum coordinate producer
identifies the whole original carrier with the union of the last $r$
coordinate hyperplanes. Thus both genuine charts give the same coordinate
model with the same selected coordinate set. Positivity puts both points
in the original carrier. The preceding whole-value transport sends each
value to that same model value at zero. Taking the composite of one
equality and the inverse of the other proves the assertion.
\end{proof}

\begin{lemma}[The first entry of the actual ambient-dimensional canonical invariant]
\label{hir:general-nc-ambient-canonical-initial-first-entry}
\leanok
\lean{Hironaka.HironakaNormalCrossingsCanonicalInitialStep.canonicalInitialValue_lt_of_branch_count_lt}
Assume $n>0$. If $\nu_H(x)<\nu_H(y)$, then
$V^{R_n}_H(x)<V^{R_n}_H(y)$ for the actual ambient-dimensional canonical
invariant $R_n$.
\uses{hir:general-nc-actual-prefix-first-entry,hir:canonical-centres}
\end{lemma}
\begin{proof}
Write $n=m+1$. By its defining positive-dimensional clause, $R_n$ is
exactly the family at depth $m+1=n$. Its first lexicographic component
is the actual nonmonomial order of the unchanged original datum at the
reflexive history. The actual original residual calculation identifies
this order with the intrinsic branch count; the differential strict
transform at the identity is the original closed carrier itself.
Consequently the asserted strict inequality follows in the first
lexicographic component, independently of every later component. This
uses the definition at depth $n$, not a comparison with depth $n+1$.
\end{proof}

\begin{theorem}[The actual ambient-dimensional initial canonical maximum]
\label{hir:general-nc-ambient-canonical-initial-full-maximum}
\leanok
\lean{Hironaka.HironakaNormalCrossingsCanonicalInitialStep.canonicalInitialMaximum_eq_highest_stratum}
Assume $n>0$ and that $r>0$ is an attained upper bound on the original
branch count. The whole maximum locus of $V^{R_n}_H$ on the unchanged
original cosupport is exactly $\{x:r\leq\nu_H(x)\}$.
\uses{hir:general-nc-generic-initial-highest-stratum-constant,hir:general-nc-ambient-canonical-initial-first-entry}
\end{theorem}
\begin{proof}
The original cosupport equals $H$. Specialize whole highest-stratum
constancy to the proved invariant $R_n$, whose dimension bound is exactly
$n$. Every point of count $r$ has the same whole value as an attained
reference point. Every remaining point of $H$ has strictly smaller count
and hence strictly smaller whole value by the first-component inequality.
Thus every highest-stratum point is a whole maximum. Conversely, a point
of smaller count cannot be a maximum because the reference point has a
strictly larger value. These two inclusions give the stated equality.
\end{proof}

\begin{theorem}[An actual first canonical step for the ambient-dimensional driver]
\label{hir:general-nc-original-ambient-canonical-first-step}
\leanok
\lean{Hironaka.HironakaNormalCrossingsCanonicalInitialStep.exists_canonicalInvariant_initial_blowup}
Assume $n>0$. From only a nonempty original NC carrier, one obtains an
attained positive highest count $r\leq n$, its whole highest stratum,
and a genuine blowup along that very stratum. The actual one-step tower
is a test tower for $M_H$ and is canonical for
$\mathrm{canonicalInvariant}(n)$ itself.
\uses{hir:general-nc-ambient-canonical-initial-full-maximum,hir:general-nc-highest-stratum-actual-blowup}
\end{theorem}
\begin{proof}
The finite range of actual branch counts, bounded by $n$, supplies an
attained positive maximum from the nonempty original carrier. The genuine
highest-stratum blowup producer gives the smooth centre and an actual
blowup. The preceding theorem identifies that same centre with the whole
maximum of the exact ambient-dimensional invariant. Its proved maximum
admissibility law therefore makes this actual blowup a test step for
the unchanged original marking-one datum. The canonical-step constructor
uses the derived whole-maximum equality and the reflexive initial history,
so it gives canonicity for the same $R_n$. No maximum, centre, history or
termination certificate is an input. No second step, arbitrary marked
datum, or identification of later branch maxima with full canonical
maxima is asserted.
\end{proof}
